\documentclass[letterpaper]{article}
\usepackage[submission]{aaai2027}
\usepackage[hyphens]{url}
\usepackage{graphicx}
\urlstyle{rm}
\def\UrlFont{\rm}
\usepackage{natbib}
\usepackage{caption}
\frenchspacing
\usepackage{amsmath,amssymb,amsfonts}
\usepackage{booktabs}
\usepackage{multirow}
\usepackage{array}
\usepackage{xspace}
\usepackage{algorithm}
\usepackage{algorithmic}
\pdfinfo{
/TemplateVersion (2027.1)
}

\setcounter{secnumdepth}{0}

\newcommand{\mars}{\textsc{MARS}\xspace}
\newcommand{\uhk}{\textsc{UHK}\xspace}
\newcommand{\mek}{\textsc{MEK}\xspace}
\newcommand{\react}{\textsc{ReAct}\xspace}
\newcommand{\codeact}{\textsc{CodeAct}\xspace}
\newcommand{\selfdebug}{\textsc{Self-Debug}\xspace}
\newcommand{\gptmini}{GPT-4o-mini\xspace}
\newcommand{\gptfour}{GPT-4o\xspace}
\newcommand{\hms}{\ensuremath{\mathrm{HMS}}\xspace}
\newcommand{\sa}{\ensuremath{\mathrm{SA}}\xspace}
\newcommand{\sr}{\ensuremath{\mathrm{SR}}\xspace}
\newcommand{\ver}{\ensuremath{\mathrm{VER}}\xspace}
\newcommand{\fileid}[1]{\texttt{\small #1}}

\title{\mars: Typed Executable Hypothesis Induction for Scientific Reasoning with Small Language Models}

\author{Anonymous Submission}
\affiliations{}

\begin{document}
\maketitle

\begin{abstract}
Scientific reasoning benchmarks expose a failure mode that is easy to miss in
standard agent evaluations: the difficult step is often not verifying a final
answer, but constructing the right hypothesis family before the answer is
known. This is especially problematic for small language models, whose extended
reasoning traces can repeat shallow variants of the same wrong representation.
We introduce \mars, a typed executable hypothesis-induction system that represents
hypotheses as executable artifacts with typed holes, evidence contracts,
measurements, validators, and residuals. A single \textsc{MARSContext} is shared
across contract compilation, typed-hole closure, estimand synthesis,
coordinate probing, metamorphic validation, and residual-conditioned operator
promotion. The model therefore proposes local typed objects, while execution
closes uncertain slots and failures are stored as structured residuals rather
than as untyped bad answers. We evaluate the same kernel on DiscoveryBench-style
data-driven discovery, NewtonBench-style symbolic law induction, and
UltraHorizon long-horizon rule induction. With a \gptmini core, \mars reaches
40.5 HMS and 60.5 consistency-HMS on a 30-task DiscoveryBench hard-bucket
slice, and 20.4 HMS / 27.1 consistency-HMS on a 239-task DB-Real audit after
residual-driven promotion of original--replication study-design operators. On
this audit, \mars exceeds published non-oracle GPT-4o ReAct and GPT-4-preview
CodeGen references, while remaining below oracle-feedback Reflexion, which
receives gold-derived revision signals. On
NewtonBench it reaches 41.7\% all-task and 83.3\%
answered-task symbolic accuracy on 24 hard-law tasks, and on UltraHorizon it
scores 49.2/45.8 on all-environment easy/hard audits while recovering 5/5 exact
programs in a checked sequence slice. A residual-operator audit further
improves sequence recovery from 2/5 to 5/5, repairs a near-correct power law
to exact recovery, and raises the most difficult DiscoveryBench raw
meta-regression subset from 0.7 to 24.2 HMS. These experiments support a
narrow structural claim:
externalizing hypothesis construction into typed executable state localizes
small-model failures and makes missing hypothesis operators measurable. Broader
operator coverage and semantic rendering remain open.
\end{abstract}

\section{Introduction}

Scientific discovery tasks ask an agent to do more than answer a question. The
agent must decide what should be measured, what kind of relation is plausible,
which variables play stable roles, and when a proposed explanation has been
falsified. Recent benchmarks make this distinction explicit. DiscoveryBench
casts data-driven discovery as hypothesis search over real datasets and
metadata \citep{discoverybench}. NewtonBench evaluates interactive scientific
law discovery rather than static curve fitting \citep{newtonbench}.
UltraHorizon stresses long-horizon partially observable exploration
\citep{ultrahorizon}, while ScienceAgentBench evaluates whether agents can
generate executable scientific workflows \citep{scienceagentbench}. Across
these settings, the bottleneck is often not final verification but the
construction of a useful hypothesis space before the answer is known.

This bottleneck is especially sharp for small language models. Standard
inference-time scaffolds extend the trajectory: \react interleaves reasoning
and actions \citep{react}, Reflexion and Self-Refine add verbal feedback
\citep{reflexion,selfrefine}, Tree-of-Thoughts and LATS search over reasoning
branches \citep{tot,lats}, and \codeact uses executable code as the action
interface \citep{codeact}. These methods improve many agent tasks, but they
still largely rely on the model to propose the intermediate concepts that make
the search fruitful. If the model never represents the task in the right
coordinates, more attempts can simply explore a larger set of wrong programs.

\mars starts from a different premise: hypothesis generation should be treated
as an external, executable induction problem, not as a longer natural-language
trace. A scientific hypothesis is represented as an artifact with typed holes,
evidence contracts, executable measurements, validation tests, and residuals.
The language model proposes local moves, but the environment closes uncertain
parts whenever possible. For example, instead of asking the model to guess a
complete law or discovery statement, the system asks which roles must be filled,
which measurement can identify each role, which residual remains after
execution, and which operator would make that residual simpler on future tasks.

The central object in \mars is a shared hypothesis state, \textsc{MARSContext}.
It is read and written by heterogeneous components: a universal hypothesis
kernel, evidence-contract compiler, estimand synthesizer, role canonicalizer,
coordinate-probe engine, metamorphic verifier, and residual compiler. This is
not a transcript memory. It is a typed state over hypotheses: what has been
measured, which slots have been closed, which answer forms are admissible,
which residuals recur, and which executable operators compress previous
failures. In this respect \mars is closer in spirit to library-learning systems
such as DreamCoder \citep{dreamcoder} than to a pure agent loop, but its learned
unit is not a domain-specific program alone. It is a reusable way of
turning an underspecified scientific task into measurable hypothesis objects.
In ReAct-style agents, a failed tool result remains contextual feedback to the
model; in \mars, it is mapped to a typed residual whose type constrains the
next admissible executable operators.

The contribution is threefold. First, we formulate small-model scientific
reasoning as \emph{typed executable hypothesis induction}: a sequence of typed
measurements, validations, and residual-conditioned language edits. Second, we
instantiate this formulation in one shared kernel that spans tabular discovery,
symbolic law induction, and long-horizon rule discovery without benchmark-
specific solver modules. Third, we report controlled small/large-model slices,
a DiscoveryBench non-oracle/oracle comparison, and a residual-operator audit
across three benchmark families. The empirical claim is deliberately narrower
than ``small models solve science'': externalizing hypothesis generation makes
failures more local and more diagnosable. The remaining gap becomes operator
coverage, full-protocol scale, and rendering alignment rather than
unconstrained chain-of-thought quality.

\section{Preliminaries}

\paragraph{Task model.}
We consider a scientific reasoning task as an interface $x$ and an answer space
$\mathcal{A}$. The interface may expose a dataset with metadata, a simulator
that can be probed, a partially observable environment, or a code workspace.
The answer may be a natural-language hypothesis, a symbolic law, a transition
program, or a Python workflow. Each benchmark provides its own evaluator
$S_x:\mathcal{A}\rightarrow\mathbb{R}$, but the system must operate before
seeing the gold answer and cannot assume a benchmark-specific answer template.

\paragraph{Hypothesis families.}
Let $\mathcal{H}$ be a family of executable hypothesis artifacts. A hypothesis
$h\in\mathcal{H}$ is not a string alone but a tuple
\begin{equation}
  h=(z,\theta,m,v,r),
\end{equation}
where $z$ is a typed sketch, $\theta$ is a partial assignment to its holes,
$m$ is an executable measurement or transition program, $v$ is a verifier, and
$r$ renders verified evidence into the benchmark answer format. This separates
three operations that are often collapsed in an LLM prompt: choosing a form,
measuring its slots, and writing the final answer.

\paragraph{Evidence contracts.}
Before proposing a final answer, \mars compiles the interface into an evidence
contract
\begin{equation}
  C(x)=\{s_i=(\tau_i,q_i,\nu_i)\}_{i=1}^{k}.
\end{equation}
Here $\tau_i$ is the slot type, $q_i$ is the observation or computation that
could close the slot, and $\nu_i$ is a local validity predicate. A contract for
tabular discovery may require population, outcome, predictor, direction, and
statistic slots. A contract for law induction may require coordinates,
dimensionless groups, constants, and residual checks. A contract for
long-horizon rule discovery may require state variables, transition clauses,
and trace-level invariants.

\paragraph{Residuals and posterior scoring.}
Executing $h$ on $x$ produces evidence $E(h,x)$ and a typed residual
\begin{equation}
  \rho(h,C)=C-\operatorname{Cover}(E(h,x),C).
\end{equation}
The kernel scores candidates by evidence fit, contract coverage, validation
loss, and complexity:
\begin{equation}
  \pi(h\mid x)\propto
  \exp\{\lambda_E E+\lambda_C\operatorname{Cover}
  -\lambda_V V-\lambda_K K\}.
\end{equation}
This resembles Bayesian model selection in spirit: a useful hypothesis should
explain observations while remaining compact and stable under checks. The
important operational point is that $\rho$ is typed. A failure is not just a low
score; it identifies a missing role, wrong coordinate system, invalid answer
form, or unmodeled transition.

Table~\ref{tab:score_terms} makes the scoring parameterization explicit. The
same decomposition is used across interfaces, while each operator normalizes
its executable evidence to a comparable range before entering the shared
posterior.

\begin{table}[t]
\centering
\small
\caption{Posterior terms used by the shared hypothesis kernel.}
\label{tab:score_terms}
\begin{tabular}{@{}llll@{}}
\toprule
Term & Definition & Range & Weight \\
\midrule
$E$ & executable fit or local score & $[0,1]$ & 24 \\
Cover & contract slots covered & $[0,1]$ & 34 \\
$V$ & metamorphic validation loss & $[0,1]$ & 32 \\
$K$ & operator complexity & $\geq0$ & 2 \\
\bottomrule
\end{tabular}
\end{table}

\paragraph{Evaluation metrics.}
DiscoveryBench reports Hypothesis Matching Score (\hms), a facet-based score
over context, variables, and relations. NewtonBench reports symbolic accuracy
(\sa), RMSLE, and interactive discovery statistics over 324 tasks spanning 12
physics domains. UltraHorizon reports environment scores for long-horizon
partially observable exploration, with trajectories designed to stress memory,
planning, and tool use. ScienceAgentBench reports program validity, execution
results, success rate, and cost over 102 tasks from 44 publications. Our main
experiments use the first three benchmark families; ScienceAgentBench remains
the natural workflow-execution target for the same hypothesis-state interface.

\section{Related Work}

\paragraph{Reasoning-time scaffolds for LLM agents.}
Chain-of-thought and agentic prompting extend the computation performed at
inference time. \react couples reasoning traces with actions
\citep{react}; Reflexion and Self-Refine use feedback to revise later attempts
\citep{reflexion,selfrefine}; Tree-of-Thoughts and LATS make the search tree
explicit \citep{tot,lats}; \codeact unifies tool use around executable Python
actions \citep{codeact}. These methods are closest to \mars at the level of
execution, but they do not by themselves define what a scientific hypothesis
object is. \mars uses the same broad agentic insight--interleave model calls
with external computation--but makes the persistent unit a typed hypothesis
state rather than a reasoning transcript.

\paragraph{Self-optimizing language-model pipelines.}
DSPy treats LM programs as optimizable computational graphs rather than
hand-written prompt strings \citep{dspy}. Self-Discover asks a model to compose
reasoning structures before solving a task \citep{selfdiscover}. Voyager stores
executable skills learned through interaction with Minecraft \citep{voyager}.
These systems show that the right unit of adaptation may be a program, module,
or reusable skill. \mars differs in two ways. First, its stored unit is tied to
scientific evidence: contracts, measurements, residuals, and validators.
Second, promotion is residual-conditioned; an operator is useful only if it
closes typed holes or compresses recurrent failures across tasks.

\paragraph{Program induction and library learning.}
DreamCoder learns a domain language and neural search policy through a
wake-sleep loop over solved and imagined programs \citep{dreamcoder}. Later
library-learning work further explores how abstractions reduce search breadth
and depth. This line is a major conceptual ancestor of \mars: both systems
treat expertise as a changing language rather than a fixed prompt. The
difference is that \mars operates in mixed scientific interfaces where the
answer may be prose, a law, or a workflow, and where holes can often be closed
by measurements rather than by full program synthesis.

\paragraph{Symbolic regression and scientific law discovery.}
Symbolic regression recovers compact formulas from data. AI-Feynman exploits
physics-inspired structure such as symmetry and separability
\citep{aifeynman}; PySR provides a practical evolutionary symbolic-regression
engine for scientific modeling \citep{pysr}. NewtonBench moves beyond static
formula fitting by requiring interactive rediscovery of shifted physical laws
\citep{newtonbench}. \mars borrows the insistence on executable hypotheses, but
does not assume that every task reduces to equation recovery. Its coordinate
probes are one operator family inside a broader hypothesis-state kernel.

\paragraph{Causal and experimental-design perspectives.}
Scientific hypotheses are not merely predictors; they make claims about stable
relations under interventions, environments, or perturbations. Structural
causal models formalize this distinction \citep{pearl}, invariant prediction
uses stability across environments to identify causal structure
\citep{invariant}, and Bayesian experimental design studies how to choose
measurements that maximize information \citep{bed}. \mars is not a causal
discovery algorithm, but it adopts the same discipline: a hypothesis must state
what evidence would support it, what transformation should preserve it, and
what residual would falsify or refine it.

\paragraph{Benchmarks for scientific agents.}
DiscoveryBench, NewtonBench, UltraHorizon, and ScienceAgentBench expose
different faces of the same problem: the agent must build an internal research
object before it can answer. DiscoveryBench emphasizes data and metadata
alignment \citep{discoverybench}; NewtonBench emphasizes interactive law
induction \citep{newtonbench}; UltraHorizon emphasizes long-horizon memory and
partial observability \citep{ultrahorizon}; ScienceAgentBench emphasizes
executable scientific workflows \citep{scienceagentbench}. We use these
benchmarks not as separate engineering targets but as probes of whether one
hypothesis-state representation can support multiple scientific task
interfaces.

\section{\mars}

\subsection{Overview}

\mars is a hypothesis-induction kernel around a small language model. The
model is used to propose typed sketches, measurements, and repairs, but the
state of the search is external and executable. The same loop is used for all
benchmarks: compile a contract, generate artifacts, execute measurements,
validate them, store residuals, and occasionally promote a residual pattern into
a new operator. Operators do not pass only text to later calls; they update one
typed context containing contracts, artifacts, evidence, residuals, and
promoted operators.

\subsection{The \textsc{MARSContext}}

For a task $x$, the context is
\begin{equation}
  \mathcal{M}_x=(C,\mathcal{H},\Theta,\mathcal{E},
  \mathcal{R},\mathcal{O},h^\star),
\end{equation}
where $C$ is the evidence contract, $\mathcal{H}$ is the pool of candidate
artifacts, $\Theta$ stores closed typed holes, $\mathcal{E}$ stores executable
evidence, $\mathcal{R}$ stores typed residuals, $\mathcal{O}$ is the active
operator set, and $h^\star$ is the current best answer artifact. All operators
read and write this object. This is the main architectural constraint: every
stage must return a state update, not a private chain of thought.

Each operator $O_j\in\mathcal{O}$ has the same interface,
\begin{equation}
  O_j(\mathcal{M}_x,x) \rightarrow
  \{(h,e,\rho,\Delta\mathcal{M})\},
\end{equation}
where $h$ is a proposed artifact, $e$ is executed evidence, $\rho$ is the typed
residual, and $\Delta\mathcal{M}$ is the state update. This interface is what
makes the system benchmark-independent. A tabular estimand, a symbolic
coordinate probe, and a transition-rule inducer differ internally, but the
kernel only sees artifact, evidence, residual, and update.

\subsection{End-to-End Trace}

Figure~\ref{fig:trace_example} gives a concrete NewtonBench-style trajectory.
The model does not guess the law directly. It proposes a typed sketch;
execution detects residual curvature; the residual selects a coordinate probe;
the probe closes the exponent and constant slots; and the held-out validator
accepts the artifact. The failed first probe changes the typed state, not
merely the next prompt.

\begin{figure}[t]
\centering
\begin{tabular}{p{0.18\linewidth}p{0.70\linewidth}}
\toprule
State & Update \\
\midrule
Contract & symbolic law; slots: coordinate, exponent, constant, residual \\
Sketch & $y=c\phi(x)^p$ with holes $\{\phi,p,c\}$ \\
Probe & raw coordinate fit leaves \texttt{wrong\_coordinate} residual \\
Closure & log-log probe sets $\phi(x)=x$, $p=2.03\rightarrow2$, $c=0.99\rightarrow1$ \\
Validate & held-out RMSLE below threshold; residual cleared; render $y=x^2$ \\
\bottomrule
\end{tabular}
\caption{A compact \mars trace. Typed residuals constrain later executable
operators: here a coordinate residual activates a log-log probe and closes the
law slots.}
\label{fig:trace_example}
\end{figure}

\subsection{Operators}

\paragraph{Contract compiler.}
The contract compiler maps the input interface to required slots and local
validators. It does not predict the gold answer. It predicts what any valid
answer must make observable: answer type, scope, variables, measurement target,
relation form, and rendering constraints.

\paragraph{Typed-hole closure.}
Instead of asking the model for a complete hypothesis, \mars asks for sketches
with typed holes and then closes the holes by small executable probes whenever
possible. For a tabular question, a hole may be an outcome role or a window
operator. For a law-discovery task, it may be a coordinate chart or exponent.
For a sequence task, it may be a state variable or transition clause. Closing a
hole adds an assignment to $\Theta$ and reduces the search space for later
holes.

\paragraph{Executable estimands and coordinate probes.}
The estimand operator creates statistical objects for tabular discovery:
group contrasts, interaction effects, coefficient shifts, multi-item
proportions, and time-window measurements. The coordinate-probe operator
creates compact transformations for laws and rules: log-log fits, affine and
polynomial lifts, ratios, finite differences, thresholds, and simple
state-transition clauses. Both operators produce evidence objects rather than
only prose rationales.

\paragraph{Metamorphic verifier.}
The verifier checks whether the artifact answers the requested question under
transformations that should preserve the object: renaming irrelevant columns,
holding answer form fixed, changing display order, perturbing nuisance values,
or replaying held-out traces. These checks are label-free and therefore
available before benchmark evaluation.

\subsection{Induction Algorithm}

\begin{algorithm}[t]
\caption{\mars typed executable hypothesis induction}
\label{alg:mars}
\begin{algorithmic}[1]
\REQUIRE task interface $x$, operator set $\mathcal{O}_0$, budget $B$
\STATE $C \leftarrow \textsc{CompileContract}(x)$
\STATE $\mathcal{M}\leftarrow(C,\emptyset,\emptyset,\emptyset,\emptyset,\mathcal{O}_0,\emptyset)$
\FOR{$t=1$ to $B$}
  \STATE $\mathcal{P}\leftarrow\emptyset$
  \FORALL{$O_j\in\mathcal{M}.\mathcal{O}$}
    \STATE $\mathcal{P}\leftarrow\mathcal{P}\cup O_j(\mathcal{M},x)$
  \ENDFOR
  \FORALL{$(h,e,\rho,\Delta\mathcal{M})\in\mathcal{P}$}
    \STATE update $\mathcal{M}$ with $h,e,\rho,\Delta\mathcal{M}$
    \STATE score $h$ by evidence, coverage, validation, and complexity
    \IF{$h$ improves $h^\star$}
       \STATE $h^\star \leftarrow h$
    \ENDIF
  \ENDFOR
  \IF{$h^\star$ satisfies $C$ and validation checks}
     \RETURN $\textsc{Render}(h^\star,C)$
  \ENDIF
  \STATE $\mathcal{O}_{new}\leftarrow\textsc{CompileResiduals}(\mathcal{M}.\mathcal{R})$
  \STATE promote operators in $\mathcal{O}_{new}$ that pass cross-task checks
\ENDFOR
\RETURN $\textsc{Render}(h^\star,C)$
\end{algorithmic}
\end{algorithm}

The score used in line 10 is the posterior-style objective from the
preliminaries:
\begin{equation}
\begin{aligned}
  \operatorname{score}(h)
  ={}& \lambda_E E(h,x)
  + \lambda_C \operatorname{cover}(h,C) \\
  &- \lambda_V V(h,C)
  - \lambda_K K(h).
\end{aligned}
\end{equation}
Here $V$ is metamorphic validation loss and $K$ is an artifact complexity
penalty. The same score ranks natural-language discovery hypotheses, symbolic
laws, and transition programs.

\subsection{Residual-Conditioned Operator Promotion}

The residual-extension step is not arbitrary code growth. A residual cluster
$R=\{\rho_1,\ldots,\rho_m\}$ can propose an operator only if it identifies a
repeatable missing transformation. Formally, a candidate operator $O'$ is
promoted when
\begin{equation}
  \Delta(O') =
  \mathbb{E}_{x\sim\mathcal{D}_{held}}
  [S(h^\star_{O'}(x))-S(h^\star(x))]
  -\beta K(O') > 0,
\end{equation}
and the improvement is not confined to the task that produced the residual.
This rule is meant to prevent an ever-growing bag of benchmark-specific tricks:
an operator must improve held-out tasks under the same contract and scoring
interface while paying a complexity penalty.

\subsection{Why This Helps Small Models}

\mars changes the role of the language model. The model does not need to hold
the entire long-horizon proof in context or invent the final hypothesis in one
sample. It repeatedly proposes local typed objects whose consequences can be
executed. Each closed hole constrains later holes, each verifier rejects an
invalid answer form before final scoring, and each residual becomes a structured
training signal for the outer loop. The intended gain is therefore not more
verbal reasoning, but a better external search geometry for weak models.

\section{Experimental Evaluation}

\subsection{Benchmarks}

We evaluate on three benchmark families that stress different interfaces but
share the same underlying difficulty: the agent must construct a hypothesis
object before it can answer. DiscoveryBench provides real datasets, metadata,
and natural-language discovery goals; the final answer is scored by
Hypothesis Matching Score (HMS) and a consistency-HMS variant. NewtonBench
provides interactive scientific-law discovery tasks in shifted physical worlds;
we report symbolic accuracy over all tasks (SA-all) and over non-abstained
answers (SA-answered). UltraHorizon provides long-horizon partially observable
exploration tasks; we report the environment-compatible score and exact
program-fit count for sequence-rule tasks. Table~\ref{tab:benchmarks} gives
the interfaces and metrics used in this draft.

\begin{table*}[t]
\centering
\caption{Benchmark interfaces used to evaluate whether one hypothesis-state
kernel can operate across heterogeneous scientific reasoning tasks.}
\label{tab:benchmarks}
\begin{tabular}{llll}
\toprule
Benchmark & Interface & Required artifact & Metric used here \\
\midrule
DiscoveryBench & tables + metadata + goal & scientific hypothesis & HMS / consistency-HMS \\
NewtonBench & simulator + observations & symbolic law & SA-all / SA-answered \\
UltraHorizon & partially observable environment & transition program & score / exact program fit \\
\bottomrule
\end{tabular}
\end{table*}

\subsection{Baselines and protocol}

The model core is either \gptmini or \gptfour. We compare raw model/agent
behavior to the same model wrapped by \mars whenever matched runs are
available. Raw baselines include direct answering and simple tool-using agent
loops; the paper treats \react, \codeact, and self-debugging as reference
families rather than as our contribution. The important controlled comparison
is whether the same model improves when hypothesis generation is routed through
contracts, executable artifacts, typed residuals, and the shared state.

We do not introduce benchmark-specific modules in the reported system. The
operator set is shared: contract compilation, typed-hole closure, estimand
synthesis, coordinate probing, metamorphic validation, and residual-conditioned
promotion. Benchmark adapters only translate the native task interface into the
generic contract representation and translate the final artifact back to the
benchmark answer format. All reported scores use each benchmark's native
metric; no cross-benchmark normalization is applied.

\begin{table*}[t]
\centering
\caption{Evaluation context. Published protocols and \mars runs are deliberately
not merged into one numeric leaderboard because their task sets and judging
surfaces differ. Published scores are discussed in text; this table records the
protocol gap.}
\label{tab:baseline_context}
\begin{tabular}{lll}
\toprule
Benchmark & Published protocol reference & \mars evaluation scope \\
\midrule
DiscoveryBench & DB-Real autonomous and oracle-feedback agents & 30-task hard slice; full DB-Real audit \\
NewtonBench & 324 interactive law-discovery tasks & 24 easy and 24 hard law slices \\
UltraHorizon & three environments with repeated long-horizon runs & all-env two-seed audit; checked sequence slice \\
\bottomrule
\end{tabular}
\end{table*}

Published DiscoveryBench reference scores include 16.3 HMS for GPT-4-preview
CodeGen, 15.4 HMS for GPT-4o ReAct, and 24.5 HMS for oracle-feedback Reflexion.
The oracle setting is separated from autonomous non-oracle agents because it
uses gold-derived feedback to revise the answer. Table~\ref{tab:db_reference}
therefore reports the non-oracle comparison and the oracle upper reference
side by side. NewtonBench reports 75.9 SA for a GPT-5 vanilla agent, 65.4 for
Gemini-2.5-Pro, and 47.8 for o4-mini. UltraHorizon reports 14.33 average score
for the best free-step LLM setting and 26.52 for humans under the paper
protocol. We cite these numbers only as context; the \mars rows below are
controlled mechanism tests, not a unified leaderboard.

\begin{table}[t]
\centering
\caption{DiscoveryBench DB-Real context. Oracle feedback uses gold-derived
revision signals, so it is not the same regime as autonomous non-oracle agents.}
\label{tab:db_reference}
\begin{tabular}{llr}
\toprule
System & Feedback & HMS \\
\midrule
GPT-4o ReAct & none & 15.4 \\
GPT-4-preview CodeGen & none & 16.3 \\
\mars, 239-task audit & none & 20.4 \\
Reflexion & oracle & 24.5 \\
\bottomrule
\end{tabular}
\end{table}

\begin{table}[t]
\centering
\caption{Evaluation results under the shared \mars operator set. Metrics are
native to each benchmark family. ``Slice'' rows are controlled mechanism tests;
``audit'' rows exercise broader task coverage without claiming a full official
leaderboard submission.}
\label{tab:main_results}
\begin{tabular}{@{}lll@{}}
\toprule
Setting & Primary & Secondary \\
\midrule
Discovery 30 slice & 40.5 HMS & 60.5 Cons. \\
Discovery 239 audit & 20.4 HMS & 27.1 Cons. \\
Newton easy & 83.3 SA-all & 95.2 SA-ans. \\
Newton hard & 41.7 SA-all & 83.3 SA-ans. \\
UH all-env easy & 49.2 score & 33.3 PI \\
UH all-env hard & 45.8 score & 33.3 PI \\
UH Seq checked & 5/5 exact & 80.0 score \\
\bottomrule
\end{tabular}
\end{table}

\subsection{Main results}

Table~\ref{tab:main_results} separates controlled slices from broader audits
for the \gptmini-core \mars system. The 30-task Discovery row tests the
hard-bucket setting used during mechanism development; the full Discovery row
is an audit over 239 DB-Real tasks. NewtonBench is reported on the
24 easy and 24 hard law-version slices. UltraHorizon is reported both as an
all-environment two-seed audit and as a checked sequence-induction slice. The
table is not a claim of full-benchmark SOTA. It tests whether the same
hypothesis-state interface transfers across three task types. The pattern is
consistent with the method design: DiscoveryBench is limited by answer-form
induction and semantic rendering, NewtonBench by coverage and abstention, and
UltraHorizon sequence induction benefits most from exact executable transition
checks. The Discovery full audit remains below the controlled slice; it
identifies residual families rather than supporting a leaderboard claim. The
post-audit original--replication operator raises the full DB-Real audit from
13.6 to 20.4 HMS, with most of the gain concentrated in meta-regression and
requirements-style table questions. In the published DB-Real context, this
places the autonomous \gptmini-core system above the non-oracle GPT-4o ReAct
and GPT-4-preview CodeGen references, but still below oracle-feedback
Reflexion.

\begin{table}[t]
\centering
\caption{Matched model comparison on identical task subsets where paired
\gptfour and \gptmini runs are available. UH rows are separate controlled
subsets and should not be compared to the 5/5 checked sequence slice in
Table~\ref{tab:main_results}.}
\label{tab:model_comparison}
\begin{tabular}{lrr}
\toprule
Setting & \gptmini & \gptfour \\
\midrule
DB stress & 64.29 & 69.84 \\
NB hard & 41.7 & 41.7 \\
UH raw & 0.0 & 6.67 \\
UH \mars & 50.0 & 43.33 \\
UH hard & 56.67 & 56.67 \\
\bottomrule
\end{tabular}
\end{table}

\subsection{Model-core comparison}

Table~\ref{tab:model_comparison} separates the language model core from the
external hypothesis system on matched subsets. DiscoveryBench still shows a
positive \gptfour gap, consistent with semantic interpretation and answer
phrasing being important. NewtonBench and UltraHorizon show a different regime:
once the hypothesis is converted into executable probes, \gptmini can match or
exceed the paired \gptfour run under the same scaffold. The UltraHorizon rows
come from paired controlled subsets, not from the checked 5/5 sequence slice;
the point of the table is model-gap analysis, not headline scoring.

\subsection{Ablations}

DiscoveryBench is the most sensitive benchmark for hypothesis quality because
the final answer is a natural-language scientific relation. Table
\ref{tab:discovery_ablation} reports the hard-bucket layer ablation used during
mechanism development. The trend is monotone: evidence contracts alone are not
enough, executable estimands improve measured hypothesis quality, role
canonicalization improves both HMS and consistency, and the final role layer
mainly increases consistency by preventing proxy columns from being treated as
unrelated scientific variables.

\begin{table}[t]
\centering
\caption{DiscoveryBench layer ablation on the hard-bucket task set.}
\label{tab:discovery_ablation}
\begin{tabular}{lrrr}
\toprule
System & N & HMS & Cons-HMS \\
\midrule
\mek & 30 & 16.7 & 30.0 \\
\mek+estimands & 30 & 20.5 & 37.2 \\
+role canonicalization & 30 & 27.5 & 44.2 \\
+universal role layers & 30 & 40.5 & 60.5 \\
\bottomrule
\end{tabular}
\end{table}

Table~\ref{tab:mechanism_ablation} summarizes additional mechanism ablations
across the three benchmark families. These are diagnostic matched-scope tests,
not full leaderboard submissions. They show the intended pattern: typed
execution helps most when the missing object is measurable. On DiscoveryBench,
residual-conditioned study-design operators repair a full-audit residual
family. On NewtonBench, coordinate charts turn a hard trigonometric slice from
mostly failed or abstained answers into exact executable laws on the matched
tasks. On UltraHorizon, typed slots and measurement compilation provide smaller
but consistent gains on the all-environment audit.

\begin{table*}[t]
\centering
\caption{Mechanism ablations across benchmark families. Each row compares
matched evaluation scopes and native metrics.}
\label{tab:mechanism_ablation}
\begin{tabular}{llll}
\toprule
Benchmark & Mechanism & Scope & Before$\to$After \\
\midrule
DiscoveryBench & hard-bucket role stack & 30 tasks & 16.7$\to$40.5 HMS \\
DiscoveryBench & residual study-design operator & 239 tasks & 13.6$\to$20.4 HMS \\
NewtonBench & self-induced coordinate charts & 4 hard tasks & 25.0$\to$75.0 SA-all \\
NewtonBench & final epistemic checks & 24 easy tasks & 79.2$\to$83.3 SA-all \\
UltraHorizon & typed slots + measurements & 6 easy env runs & 46.7$\to$50.8 score \\
UltraHorizon & sequence residual operators & checked hard seq & 2/5$\to$5/5 exact \\
\bottomrule
\end{tabular}
\end{table*}

\subsection{Residual-operator audit}

The strongest novelty claim in \mars is the residual-to-operator pathway. We
therefore report it separately from the main score table. Table
\ref{tab:operator_audit} summarizes smoke and audit runs in which residual
patterns were converted into executable rules or promoted operator families.
These runs are not full official benchmark evaluations; they test whether
residuals can change the hypothesis language in measurable ways.

\begin{table*}[t]
\centering
\caption{Residual-operator audit. Source and promotion columns distinguish
where the residual was observed from where the candidate operator was checked.}
\label{tab:operator_audit}
\begin{tabular}{lllll}
\toprule
Slice & Operator & Source & Promotion eval. & Before$\to$After \\
\midrule
UH Seq & string rules & train traces & held-out rules & 2/5$\to$5/5 \\
Newton smoke & power law & residual fit & held-out points & 0.85$\to$1.00 \\
Discovery & study-design roles & failed tasks & DB-Real full audit & 13.6$\to$20.4 \\
Discovery raw meta-reg. & paired-slot renderer & failed tasks & 50-task subset & 0.7$\to$24.2 \\
Mean norm. & G0$\to$G1 & outer loop & suite audit & .789$\to$.801 \\
\bottomrule
\end{tabular}
\end{table*}

The audit is deliberately separated from the benchmark table because the
evidence is mechanistic rather than comprehensive. In UltraHorizon sequence
tasks, residual rules such as character-wise combination, positional maximum,
and sorted multiset transformations close failures that a raw proposal loop
missed. In the Newton smoke setting, a numeric residual kernel repairs an
almost-correct power law on held-out points, not a separate full NewtonBench
task split. In the DiscoveryBench outer-loop audit, promoted study-design
operators infer original/replication arms, domain-code aliases, paired columns,
binary complements, and contrastive renderings from table schemas. They move the
full DB-Real audit from 13.6 to 20.4 HMS and the raw meta-regression subset from
0.7 to 24.2 HMS. The remaining plateau is equally important: the current
mechanism can stabilize a better operator family, but does not yet guarantee
continuing open-ended improvement across unrelated table families.

\subsection{Failure analysis}

The remaining errors fall into four recurring classes. \emph{Rendering errors}
occur when the system measures a valid relation but expresses it in a form that
does not match the benchmark facet. These dominate DiscoveryBench. \emph{Role
errors} occur when the contract closes a plausible but wrong scientific role,
often because several dataset columns proxy the same latent variable.
\emph{Coverage errors} occur when no active operator proposes the right
coordinate chart, transition clause, or estimand family. These dominate
NewtonBench abstentions. \emph{Assembly errors} occur when the task requires a
large multi-file workflow; this is the main reason ScienceAgentBench is treated
as a target interface rather than as a headline result in this draft.

The 239-task Discovery audit makes these errors concrete. After residual
operator promotion, the system is strong when the task is close to executable
slot closure or paired study-design measurement: 73.6 HMS on the NLS raw subset,
36.0 HMS on WorldBank indicator tasks, 24.5 HMS on meta-regression, 24.2 HMS on
raw meta-regression, 26.7 HMS on requirements-engineering questions, and
23.7 HMS on archaeology tasks. The same audit exposes negative transfer and
uncovered families: plant introduction pathways fall to 6.3 HMS, NLS SES falls
to 7.4 HMS, and one WorldBank education-GDP subset remains at 0.0 HMS. The
dominant remaining DiscoveryBench failures are denominator alignment,
multi-file assembly, and semantic rendering, not arithmetic measurement.

This failure analysis is also the practical value of the \mars state. In a
transcript-only agent, failures are often just bad final answers. In \mars, the
failure is attached to the missing typed object: a role, a coordinate, a
validator, a renderer, or an operator. That makes later improvement measurable
without changing the benchmark-specific code path.

\section{Limitations}

\mars is a hypothesis-induction architecture, not evidence that small language
models can autonomously solve scientific discovery in general. The results in
this draft are controlled slices, not full official benchmark submissions.
They are useful because they expose the mechanism under matched conditions, but
they should not be read as a completed state-of-the-art claim.

The strongest results occur when the task exposes measurements or traces that
can close typed holes. When the needed operator is absent, the kernel cannot
recover the right hypothesis by ranking alone. This makes operator coverage the
main remaining technical bottleneck. DiscoveryBench additionally exposes a
rendering problem: the system can measure a scientifically plausible relation
but express it in a form that misses the benchmark's target facet. NewtonBench
errors are concentrated in abstentions, compile failures, and missing
coordinate charts rather than in confidently accepted wrong laws.
UltraHorizon sequence induction is strong in the checked setting, but the full
benchmark requires more environments and repeated runs because single
trajectories have high variance. ScienceAgentBench is excluded from the
headline evaluation because full scientific workflow execution requires a
stronger multi-file assembler.

The residual-operator audit is also limited. It demonstrates that residuals can
produce reusable operators on held-out slices, but it does not yet prove
unbounded autonomous growth. The current system has a plateau after the first
outer-loop improvement. A stronger version must learn broader operator
families from residuals while retaining the same contract, validation, and
complexity controls.

\paragraph{DiscoveryBench positioning.}
The correct reading of the DiscoveryBench result is not ``full SOTA.'' The
strong claim is that a small-model system without gold feedback can exceed the
published non-oracle GPT-4o ReAct and GPT-4-preview CodeGen references on the
DB-Real audit, while oracle-feedback Reflexion remains a higher revision-aided
reference. This distinction matters because the proposed method targets
autonomous hypothesis-state construction rather than access to answer-level
feedback.

\section{Conclusion}

We presented \mars, a typed executable hypothesis-induction system for
scientific reasoning with small language models. The main claim is structural:
weak models benefit when hypothesis generation is moved from free-form prose
into an external state of typed artifacts, executable measurements, validators,
and residuals. Across data-driven discovery, symbolic law induction, and
long-horizon sequence induction, the same kernel improves the effective
reasoning of \gptmini on controlled slices and sometimes reduces the gap to
\gptfour. The experiments do not show that scaffolding eliminates model
capability gaps or that the system has already reached full-benchmark SOTA.
They show where external structure helps: when the system can decompose a
hypothesis into measurable slots, use execution to close them, and turn
recurrent typed residuals into better operators. This suggests a concrete path
for small-model scientific agents: grow the hypothesis language itself, guided
by residuals, rather than only sampling more reasoning traces from a fixed
language.

\begin{thebibliography}{99}

\bibitem[Majumder et al.(2024)]{discoverybench}
Bodhisattwa Prasad Majumder, Harshit Surana, Dhruv Agarwal, Bhavana
Dalvi Mishra, Abhijeetsingh Meena, Aryan Prakhar, Tirth Vora, Tushar
Khot, Ashish Sabharwal, and Peter Clark.
\newblock DiscoveryBench: Towards data-driven discovery with large language
models.
\newblock arXiv preprint arXiv:2407.01725, 2024.

\bibitem[Zheng et al.(2025)]{newtonbench}
Tianshi Zheng, Kelvin Kiu-Wai Tam, Newt Hue-Nam K. Nguyen, Baixuan Xu,
Zhaowei Wang, Jiayang Cheng, Hong Ting Tsang, Weiqi Wang, Jiaxin Bai,
Tianqing Fang, Yangqiu Song, Ginny Y. Wong, and Simon See.
\newblock NewtonBench: Benchmarking generalizable scientific law discovery in
LLM agents.
\newblock arXiv preprint arXiv:2510.07172, 2025.

\bibitem[Luo et al.(2025)]{ultrahorizon}
Haotian Luo, Huaisong Zhang, Xuelin Zhang, Haoyu Wang, Zeyu Qin, Wenjie
Lu, Guozheng Ma, Haiying He, Yingsha Xie, Qiyang Zhou, Zixuan Hu,
Hongze Mi, Yibo Wang, Naiqiang Tan, Hong Chen, Yi R. Fung, Chun Yuan,
and Li Shen.
\newblock UltraHorizon: Benchmarking agent capabilities in ultra long-horizon
scenarios.
\newblock arXiv preprint arXiv:2509.21766, 2025.

\bibitem[Chen et al.(2025)]{scienceagentbench}
Ziru Chen, Shijie Chen, Yuting Ning, Qianheng Zhang, Boshi Wang,
Botao Yu, Yifei Li, Zeyi Liao, Chen Wei, Zitong Lu, Vishal Dey,
Mingyi Xue, Frazier N. Baker, Benjamin Burns, Daniel Adu-Ampratwum,
Xuhui Huang, Xia Ning, Song Gao, Yu Su, and Huan Sun.
\newblock ScienceAgentBench: Toward rigorous assessment of language agents for
data-driven scientific discovery.
\newblock In \emph{ICLR}, 2025.

\bibitem[Yao et al.(2023a)]{react}
Shunyu Yao, Jeffrey Zhao, Dian Yu, Nan Du, Izhak Shafran, Karthik
Narasimhan, and Yuan Cao.
\newblock ReAct: Synergizing reasoning and acting in language models.
\newblock In \emph{ICLR}, 2023.

\bibitem[Shinn et al.(2023)]{reflexion}
Noah Shinn, Federico Cassano, Edward Berman, Ashwin Gopinath, Karthik
Narasimhan, and Shunyu Yao.
\newblock Reflexion: Language agents with verbal reinforcement learning.
\newblock In \emph{NeurIPS}, 2023.

\bibitem[Madaan et al.(2023)]{selfrefine}
Aman Madaan, Niket Tandon, Prakhar Gupta, Skyler Hallinan, Luyu Gao,
Sarah Wiegreffe, Uri Alon, Nouha Dziri, Shrimai Prabhumoye, Yiming Yang,
Shashank Gupta, Bodhisattwa Prasad Majumder, Katherine Hermann, Sean
Welleck, Amir Yazdanbakhsh, and Peter Clark.
\newblock Self-Refine: Iterative refinement with self-feedback.
\newblock In \emph{NeurIPS}, 2023.

\bibitem[Yao et al.(2023b)]{tot}
Shunyu Yao, Dian Yu, Jeffrey Zhao, Izhak Shafran, Thomas L. Griffiths,
Yuan Cao, and Karthik Narasimhan.
\newblock Tree of thoughts: Deliberate problem solving with large language
models.
\newblock In \emph{NeurIPS}, 2023.

\bibitem[Zhou et al.(2024)]{lats}
Andy Zhou, Kai Yan, Michal Shlapentokh-Rothman, Haohan Wang, and
Yu-Xiong Wang.
\newblock Language agent tree search unifies reasoning, acting, and planning in
language models.
\newblock In \emph{ICML}, 2024.

\bibitem[Wang et al.(2024)]{codeact}
Xingyao Wang, Yangyi Chen, Lifan Yuan, Yizhe Zhang, Yunzhu Li, Hao Peng,
and Heng Ji.
\newblock Executable code actions elicit better LLM agents.
\newblock In \emph{ICML}, 2024.

\bibitem[Khattab et al.(2023)]{dspy}
Omar Khattab, Arnav Singhvi, Paridhi Maheshwari, Zhiyuan Zhang, Keshav
Santhanam, Sri Vardhamanan, Saiful Haq, Ashutosh Sharma, Thomas T.
Joshi, Hanna Moazam, Heather Miller, Matei Zaharia, and Christopher Potts.
\newblock DSPy: Compiling declarative language model calls into
self-improving pipelines.
\newblock arXiv preprint arXiv:2310.03714, 2023.

\bibitem[Zhou et al.(2024)]{selfdiscover}
Pei Zhou, Jay Pujara, Xiang Ren, Xinyun Chen, Heng-Tze Cheng, Quoc V.
Le, Ed H. Chi, Denny Zhou, Swaroop Mishra, and Huaixiu Steven Zheng.
\newblock Self-Discover: Large language models self-compose reasoning
structures.
\newblock arXiv preprint arXiv:2402.03620, 2024.

\bibitem[Wang et al.(2023)]{voyager}
Guanzhi Wang, Yuqi Xie, Yunfan Jiang, Ajay Mandlekar, Chaowei Xiao,
Yuke Zhu, Linxi Fan, and Anima Anandkumar.
\newblock Voyager: An open-ended embodied agent with large language models.
\newblock arXiv preprint arXiv:2305.16291, 2023.

\bibitem[Ellis et al.(2021)]{dreamcoder}
Kevin Ellis, Catherine Wong, Maxwell Nye, Mathias Sable-Meyer, Luc Cary,
Lucas Morales, Luke Hewitt, Armando Solar-Lezama, and Joshua B. Tenenbaum.
\newblock DreamCoder: Growing generalizable, interpretable knowledge with
wake-sleep Bayesian program learning.
\newblock In \emph{PLDI}, 2021.

\bibitem[Udrescu and Tegmark(2020)]{aifeynman}
Silviu-Marian Udrescu and Max Tegmark.
\newblock AI Feynman: A physics-inspired method for symbolic regression.
\newblock \emph{Science Advances}, 6(16), 2020.

\bibitem[Cranmer(2023)]{pysr}
Miles Cranmer.
\newblock Interpretable machine learning for science with PySR and
SymbolicRegression.jl.
\newblock arXiv preprint arXiv:2305.01582, 2023.

\bibitem[Pearl(2009)]{pearl}
Judea Pearl.
\newblock \emph{Causality: Models, Reasoning, and Inference}.
\newblock Cambridge University Press, 2nd edition, 2009.

\bibitem[Peters, Buhlmann, and Meinshausen(2016)]{invariant}
Jonas Peters, Peter Buhlmann, and Nicolai Meinshausen.
\newblock Causal inference by using invariant prediction: Identification and
confidence intervals.
\newblock \emph{Journal of the Royal Statistical Society: Series B},
78(5):947--1012, 2016.

\bibitem[Rainforth et al.(2023)]{bed}
Tom Rainforth, Adam Foster, Desi R. Ivanova, and Freddie Bickford Smith.
\newblock Modern Bayesian experimental design.
\newblock arXiv preprint arXiv:2302.14545, 2023.

\end{thebibliography}

\end{document}
